%!TEX root = ../template.tex
%%%%%%%%%%%%%%%%%%%%%%%%%%%%%%%%%%%%%%%%%%%%%%%%%%%%%%%%%%%%%%%%%%%%
%% appendix-triggers.tex
%% NOVA thesis document file
%%
%% Appendix: why the timing gate fired (FGSM-parity runs), in detail
%%%%%%%%%%%%%%%%%%%%%%%%%%%%%%%%%%%%%%%%%%%%%%%%%%%%%%%%%%%%%%%%%%%%

\typeout{NT FILE appendix-triggers.tex}%

\chapter{Timing Gate Triggers}
\label{app:triggers}

This appendix gives the full trigger breakdown of the FGSM-parity evaluations, whose
conclusions are stated in Sections~\ref{subsec:results_parity_triggers}
and~\ref{subsec:results_parity_coord}. These runs are the first in which the timing gate
recorded why it acted (Section~\ref{subsec:procedure_records}). For every attacked step
it logs which case of Equation~\eqref{eq:timing_rule} admitted the step, a score
strictly above the threshold or a tie at the threshold admitted by rationing, together
with the congestion band the score fell in: saturated ($c_t \ge 1.00$), near saturation
($0.90 \le c_t < 1.00$), high congestion ($0.75 \le c_t < 0.90$) or moderate
($c_t < 0.75$). Figures~\ref{fig:trigger_reasons_parity}
and~\ref{fig:trigger_reasons_coordinated_parity} give each trigger's share of the attack
arm's attacked steps in the two scopes. Ungated runs have no gate and are omitted.

% ─────────────────────────────────────────────────────────────────────────────
\section{Independent scope}
\label{app:triggers_independent}

\begin{figure}[htbp]
  \centering
  \includegraphics[width=\textwidth]{trigger_reasons_independent_parity.png}
  \caption{Independent scope, FGSM-parity domain clamp: why the timing gate fired in
           the attack arm ($\varepsilon = 0.30$, $n = 30$ episodes). One panel per
           trigger that occurred in at least one run, combining the case of
           Equation~\eqref{eq:timing_rule} that admitted the step with the
           congestion band of its score. Rows are victim variants and columns the
           requested timing budget. Each cell gives the trigger's share of that run's
           attacked steps, with the count beneath it, so the three panels of one cell
           sum to 100\%. Ungated runs have no gate and are omitted.}
  \label{fig:trigger_reasons_parity}
\end{figure}

No attacked step in any run, in either arm, scored $1.00$ or more, or less than $0.75$.
At $b = 0.10$ and $b = 0.15$ the mean threshold at which the gate fired is $0.90$ for
every variant, and the tied steps are the steps sitting on the point mass at that value.
Its stored value falls a floating-point rounding error below $0.90$, which is why ties
are recorded in the high-congestion band although they sit at that band's upper edge.

Steps scoring between $0.90$ and $1.00$ are admitted at every budget, in a nearly
constant number per variant: $321$ for \texttt{CC-Simple} and between $505$ and $599$
for the others. Their share of the attacks therefore falls from $42$--$77\%$ at
$b = 0.10$ to $17$--$32\%$ at $b = 0.25$.

The route by which the rest of the budget is spent changes with the budget. At
$b = 0.10$ and $b = 0.15$ the threshold sits on the point mass, and ties admitted by
rationing account for $23$--$58\%$ and $45$--$66\%$ of the attacks respectively. On the
observed score sequences, admitting every tie would have raised the attack rate to
between $17\%$ and $21\%$ at both budgets. From $b = 0.20$ the threshold begins to fall
below the point mass, and by $b = 0.25$ it lies between $0.82$ and $0.90$. The steps on
the point mass then pass as strictly above the threshold, and above-threshold steps in
the high-congestion band, the point mass among them, make up $57$--$77\%$ of the attacks
at $b = 0.25$, while ties fall to $6$--$11\%$.

The random control passes through the same gate, and its trigger composition is nearly
identical: at $b = 0.15$, for instance, its tie share differs from the attack arm's by
at most three percentage points for every variant.

% ─────────────────────────────────────────────────────────────────────────────
\section{Coordinated scope}
\label{app:triggers_coordinated}

\begin{figure}[htbp]
  \centering
  \includegraphics[width=\textwidth]{trigger_reasons_coordinated_parity.png}
  \caption{Coordinated scope, FGSM-parity domain clamp: why the timing gate fired in the
           attack arm ($\varepsilon = 0.30$, $n = 30$ episodes). Panels, rows, columns
           and cell contents as in Figure~\ref{fig:trigger_reasons_parity}.}
  \label{fig:trigger_reasons_coordinated_parity}
\end{figure}

The picture is the same when the perturbation is coordinated. No attacked step scored
$1.00$ or more, or less than $0.75$. Steps scoring between $0.90$ and $1.00$ fire in a
nearly constant number per variant at every budget, so their share of the attacks falls
from $42$--$78\%$ at $b = 0.10$ to $17$--$32\%$ at $b = 0.25$. Ties admitted by
rationing make up $22$--$58\%$ of the attacks at $b = 0.10$ and $44$--$66\%$ at
$b = 0.15$; admitting every tie would have raised the attack rate to between $17.5\%$
and $21.3\%$ at both budgets. At $b = 0.25$ the threshold lies between $0.82$ and
$0.90$, above-threshold steps in the high-congestion band make up $58$--$77\%$ of the
attacks, and ties only $3$--$11\%$. Every one of these figures is within a few
percentage points of its counterpart in the independent scope, which is what the gate's
score being a function of link utilisation alone implies.
